% [原文 p.3]
理论的行为就会改变——在尺度改变之前，它可能被认作另一个不同理论的行为。理论随标度改变而发生的这种
变化称为\emph{重整化群流}，简称 \emph{RG 流}。若一个理论在 RG 流下不变，则称它是\emph{标度不变的}。
这种理论中算子 $\cO$ 的两点关联函数按幂律依赖于距离，
\begin{equation}
  \langle\cO(x)\cO(y)\rangle=\frac{1}{\mathrm{dist}(x,y)^{2\Delta_{\cO}}}. \tag{2.7}
\end{equation}
数 $\Delta_{\cO}$ 称为该算子的\emph{维数}。

在二维情形，可以证明 QFT 的标度不变性等价于共形不变性。对一般 QFT，我们在紫外与红外极限下
分别得到共形不变的场论（CFT）。CFT 的一个不变量是它的中心荷 $c$；众所周知 $c$ 在 RG 流下减小：
$c_{\mathrm{UV}}\ge c_{\mathrm{IR}}$。

同样的结论当然也适用于带超对称的 QFT。在二维 $(2,2)$ 超对称 QFT 中，存在一类在 RG 流下不变
的可观测量，即便理论本身并非标度不变。上面引入的手征环与膜的范畴便是这类“受保护”而免于重整化的
可观测量的例子。在 $(2,2)$ 超共形场论（SCFT）中，用 $\hat c=c/3$ 表示中心荷比较方便。

\subsection{形变}
\label{sec:deform}

一个 QFT 可以通过向其拉氏量密度中加入局部算子 $\cO$ 来形变。若理论是标度不变的，则当算子维数
减去时空维数（在二维理论中即 $\Delta_{\cO}-2$）为正、零、负时，分别称该形变为\emph{无关的}、
\emph{边缘的}、\emph{相关的}。在 RG 流下，无关形变衰减，相关形变增长。若一个边缘形变在 RG 流下
保持不变，则称它是\emph{精确边缘的}；若它在 RG 流下衰减（相应地增长），则称它是\emph{边缘无关的}
（相应地\emph{边缘相关的}）。共形场论的模空间由精确边缘形变参数来坐标化，它带有自然的度规 $G$，
称为 Zamolodchikov 度规 [2]：设 $v_1,v_2$ 是某个理论处对应于精确边缘算子 $\cO_1,\cO_2$ 的切向量，
则它们的内积 $G(v_1,v_2)$ 由（二维情形下）下式给出：
\begin{equation}
  \langle\cO_1(x)\cO_2(y)\rangle=\frac{G(v_1,v_2)}{\mathrm{dist}(x,y)^{4}}. \tag{2.8}
\end{equation}

% [原文 p.4]
在二维 $(2,2)$ 超对称 QFT 中，保持超对称的形变算子有如下三种类型：
\begin{equation}
  \Delta_D\cL=Q_+Q_-\overline Q_-\overline Q_+\cK, \tag{2.9}
\end{equation}
\begin{equation}
  \Delta_A\cL=Q_+\overline Q_-\cO_A\ \ \text{及其伴随}, \tag{2.10}
\end{equation}
\begin{equation}
  \Delta_B\cL=Q_+Q_-\cO_B\ \ \text{及其伴随}, \tag{2.11}
\end{equation}
其中 $\cK$ 是标量算子且任意，$\cO_A$ 是 \emph{A-手征的}，$\overline Q_+ \cO_A=Q_-\cO_A=0$；
$\cO_B$ 是 \emph{B-手征的}，$\overline Q_\pm \cO_B=0$。我们分别称它们为 \emph{D-项}、\emph{A-项}、
\emph{B-项}。\footnote{标准的术语是：用\emph{扭 F-项}（twisted $F$-term）代替 A-项，用
\emph{F-项}（$F$-term）代替 B-项。}%
由 (2.1)--(2.2) 可以看出，差一个全导数后，一个 D-项同时是一个 A-项和一个 B-项。可以证明，
A-项形变模去 D-项形变，与手征环 $\cR_A$ 的元素一一对应；类似地，B-项形变模去 D-项形变，
与手征环 $\cR_B$ 的元素一一对应。

由代数 (2.1)--(2.2) 可推出：差一个全导数后，D-项与 A-项都是 $Q_B$-恰当的，而 D-项与 B-项都是
$Q_A$-恰当的。特别地，环 $\cR_B$ 与范畴 $\cC_B$ 在 D-项形变与 A-项形变下不变，而 $\cR_A$ 与
范畴 $\cC_A$ 在 D-项形变与 B-项形变下不变。

在 $(2,2)$ SCFT 中，D-项形变是无关的，因为每个超荷都带有维数 $\frac12$。只有一部分 A-项与
B-项形变是边缘的或相关的。

对应于 A-项形变与 B-项形变的理论参数空间分别记作 $\fM_A$ 与 $\fM_B$。它们具有复结构：
$\fM_A$（相应地 $\fM_B$）的 $(1,0)$ 型切向量对应于形如 $Q_+\overline Q_-\cO_A$（相应地
$Q_+Q_-\cO_B$）的算子，并可自然地等同于手征环 $\cR_A$（相应地 $\cR_B$）的元素。
对一个 $(2,2)$ SCFT，精确边缘参数的子空间 $\fM_A^0\subset\fM_A$ 与 $\fM_B^0\subset\fM_B$
是复子流形，它们也是共形场论模空间的子流形。众所周知，诱导到 $\fM_A^0$ 与 $\fM_B^0$ 上的
Zamolodchikov 度规是 Kähler 的 [3]。

